\documentclass{article}
\usepackage[T1]{fontenc}
\usepackage{lmodern}
\usepackage{graphicx}
\usepackage{amsmath,amssymb}
\usepackage[a4paper,margin=2cm]{geometry}
\usepackage[absolute,overlay]{textpos}
\setlength{\TPHorizModule}{1mm}
\setlength{\TPVertModule}{1mm}
\usepackage[indonesian]{babel}
\usepackage{hyperref}
\setlength{\emergencystretch}{2em}
% unit: Halaman 85

\begin{document}

% === Halaman 85 (python/native) ===

Step 5: State hasil 10 putaran-ganda dijumlahkan dengan state awal yang disimpan 

pada step 3: 

              y = (x′ + x) mod 2³²

Jika plainteks berukuran n x 512 bit, maka Salsa20 dieksekusi sebanyak n kali, 

dengan counter bertambah 1, sedangkan nonce tetap.

85

Kode program C++:

\#include <stdint.h>

\#define ROTL(a,b) (((a) << (b)) | ((a) >> (32 - (b))))

\#define QR(a, b, c, d)(  \textbackslash{}


b \textasciicircum{}= ROTL(a + d, 7), \textbackslash{}


c \textasciicircum{}= ROTL(b + a, 9), \textbackslash{}


d \textasciicircum{}= ROTL(c + b,13), \textbackslash{}


a \textasciicircum{}= ROTL(d + c,18))

\#define ROUNDS 20

\end{document}
